\section{高效视频理解：采样、选择和轻量化}

\subsection{为什么效率问题会变成主问题}

长视频很长，3D 卷积很贵，attention 更贵。如果模型每次都把所有帧完整跑一遍，成本会迅速失控。因此，视频理解越来越像一个“在预算内找信息”的问题，而不是纯粹的分类问题。

\begin{knowledgebox}{效率优化的思路}
高效视频理解不是简单“把模型变小”，而是把计算集中在最有用的时间段、空间区域和模态上。它本质上是在做信息预算分配。
\end{knowledgebox}

\subsection{几类常见策略}

效率提升通常有几种思路：

\begin{enumerate}
    \item \textbf{采样更少的 clip}：只处理更有信息的片段。
    \item \textbf{选择更轻的骨干}：例如更轻量的 3D 网络或更高效的 backbone。
    \item \textbf{做模态选择}：有些任务不一定需要所有模态同时推理。
    \item \textbf{做预览机制}：先用便宜模态估计哪里值得看，再做精细处理。
\end{enumerate}

\begin{figure}[H]
\centering
\includegraphics[width=0.95\textwidth]{slides-images/slide-048.jpg}
\caption{长视频动作分类的最终效果：运动信息、外观信息和融合策略共同决定性能\protect\footnotemark}
\end{figure}
\footnotetext{Slides 第49页。}

\subsection{从长视频到 egocentric video}

更进一步，课程把话题推进到 egocentric / exocentric 结合的场景。第一视角视频和外部视角视频可以互相补充，尤其在多人与多摄像头场景里，视频理解开始变成“谁在看什么、谁在和谁交互”的问题。

\begin{warningbox}{第一视角视频更难}
第一视角视频更接近真实体验，但同时也更抖、更碎、更依赖上下文，还会把“操作者意图”引入任务。这让模型更接近人，也更难。
\end{warningbox}

